\documentclass[10pt,a4paper]{amsart}
\usepackage[margin=19mm]{geometry}
\usepackage{amsmath,amssymb,mathtools}
\usepackage[scr=rsfs,cal=euler]{mathalfa}
\usepackage{xfrac}
\usepackage[unicode,hidelinks,bookmarks=false]{hyperref}
\newcommand{\ie}{i.e.\ }
\newcommand{\cf}{cf.\ }
\newcommand{\resp}{respectively\ }
\newcommand{\Subsection}{\subsection}
\providecommand{\nobreakdash}{\nobreak\hbox{-}}
\setlength{\emergencystretch}{2em}
\multlinegap0pt
\usepackage{xcolor}
\usepackage{amssymb}
\usepackage{mathtools}
\usepackage{algorithm}\usepackage{algpseudocode}
\newtheorem{lemma}{Remark}
\title{Analysis of Shuffling Beyond Pure Local Differential Privacy}
\author{Shun Takagi, Seng Pei Liew}
\date{Source: arXiv:2601.19154v4 (CC BY 4.0)}
\begin{document}
\maketitle

\noindent\emph{Source and licence.} Analysis of Shuffling Beyond Pure Local Differential Privacy. Version arXiv:2601.19154v4 (CC BY 4.0), distributed by arXiv under the Creative Commons Attribution 4.0 International licence, http://creativecommons.org/licenses/by/4.0/, as recorded in the arXiv licence metadata for this exact version and shown on the abstract page https://arxiv.org/abs/2601.19154v4. Full licence notice: \texttt{licenses/arxiv-2601.19154-CC-BY-4.0.txt}.\par\smallskip

\noindent\textit{Editorial scope.} Counted unit: the article's lemma lem:pairwise-chi-structure (source0.tex lines 4091--4127). The article's complete original proof of it is delivered with it (source0.tex lines 4129--4283, one \texttt{proof} environment of 155 lines). No mathematics is written, repaired or re-derived: the statement and the proof are the article's own lines, in the article's own notation, and no retained line is substituted or re-flowed.  No further source-local context is needed: every cross-reference of every retained line resolves inside the two delivered spans.  Nothing was omitted from any retained span, so the package carries no omission record.  No retained line contains a citation, so the wrapper carries no bibliography and no bibliography entry.  No retained line contains a figure environment, an image-inclusion command or drawing code, so no image is delivered and the package declares no asset dependency.  Editorial formatting only, in addition: an AMS \texttt{amsart} preamble that restates the article's own declarations and macros used by the retained lines, namely the environments \texttt{lemma} on separate counters, together with the standard packages those lines need.  The retained environments print in this document as Lemma 1 in order of appearance, whereas the article prints the same items with the numbering of its own full document.  The complete native source of the article as distributed in the arXiv e-print for this exact version is bundled unchanged as \texttt{sources/193440/source0.tex}.
\begin{lemma}[Pairwise ordering and equality condition]
\label{lem:pairwise-chi-structure}
Let $x_1\neq x_1'\in\mathcal X$ be fixed and consider the pairwise
shuffle indices
\[
  \chi_{\mathrm{lo}}(x_1,x_1')
  := \sqrt{
        \frac{\gamma}{
          \mathrm{Var}_{Y\sim\mathcal R_{\mathrm{BG}}}
          \bigl(l_0(Y;x_1,x_1',\mathcal R_{\mathrm{BG}})\bigr)
        }
      },
  \qquad
  \chi_{\mathrm{up}}(x_1,x_1')
  := \inf_{x\in\mathcal X}
      \sqrt{
        \frac{1}{
          \mathrm{Var}_{Y\sim\mathcal R_x}
          \bigl(l_0(Y;x_1,x_1',\mathcal R_x)\bigr)
        }
      }.
\]
Then
\[
  \chi_{\mathrm{up}}(x_1,x_1')
  \;\ge\;
  \chi_{\mathrm{lo}}(x_1,x_1')
  \qquad\text{for all }x_1\neq x_1'.
\]
Moreover, equality holds if and only if there exists an input
$x^\ast\in\mathcal X$ such that
\[
  \mathcal R_{x^\ast}(y)
  = \underline{\mathcal R}(y)
  \qquad\text{for almost every }y\in A(x_1,x_1').
\]
\end{lemma}

\begin{proof}
Fix $x_1\neq x_1'\in\mathcal X$ and abbreviate
\[
  \Delta(y)
  := \mathcal R_{x_1}(y) - \mathcal R_{x_1'}(y),
  \qquad y\in\mathcal Y.
\]
Note that $\int_{\mathcal Y}\Delta(y)\,dy=0$, since both
$\mathcal R_{x_1}$ and $\mathcal R_{x_1'}$ are probability distributions
with respect to the same reference measure.

\textbf{Step 1: Variance under $\mathcal R_{\mathrm{BG}}$.}
For $Y\sim\mathcal R_{\mathrm{BG}}$ we have
\[
  l_0(Y;x_1,x_1',\mathcal R_{\mathrm{BG}})
  = \frac{\mathcal R_{x_1}(Y)-\mathcal R_{x_1'}(Y)}
         {\mathcal R_{\mathrm{BG}}(Y)}
  = \frac{\Delta(Y)}{\mathcal R_{\mathrm{BG}}(Y)}.
\]
Since
$\mathbb E_{\mathcal R_{\mathrm{BG}}}[l_0(Y;x_1,x_1',\mathcal R_{\mathrm{BG}})]
= 0$, its variance is
\begin{align*}
  \sigma_{\mathrm{BG}}^2(x_1,x_1')
  &:= \mathrm{Var}_{Y\sim\mathcal R_{\mathrm{BG}}}
      \bigl(l_0(Y;x_1,x_1',\mathcal R_{\mathrm{BG}})\bigr)\\
  &= \int_{\mathcal Y}
       \frac{\Delta(y)^2}{\mathcal R_{\mathrm{BG}}(y)}\,dy.
\end{align*}
On the complement of $A(x_1,x_1')$ we have $\Delta(y)=0$, so
\[
  \sigma_{\mathrm{BG}}^2(x_1,x_1')
  = \int_{A(x_1,x_1')}
      \frac{\Delta(y)^2}{\mathcal R_{\mathrm{BG}}(y)}\,dy.
\]

\textbf{Step 2: Variance under $\mathcal R_x$.}
Similarly, for any $x\in\mathcal X$ and $Y\sim\mathcal R_x$,
\[
  l_0(Y;x_1,x_1',\mathcal R_x)
  = \frac{\mathcal R_{x_1}(Y)-\mathcal R_{x_1'}(Y)}
         {\mathcal R_x(Y)}
  = \frac{\Delta(Y)}{\mathcal R_x(Y)},
\]
and again the mean is zero. Hence
\begin{align*}
  \sigma_x^2(x_1,x_1')
  &:= \mathrm{Var}_{Y\sim\mathcal R_x}
      \bigl(l_0(Y;x_1,x_1',\mathcal R_x)\bigr)\\
  &= \int_{\mathcal Y}
       \frac{\Delta(y)^2}{\mathcal R_x(y)}\,dy
   = \int_{A(x_1,x_1')}
       \frac{\Delta(y)^2}{\mathcal R_x(y)}\,dy.
\end{align*}

\textbf{Step 3: Comparison of variances.}
By definition of the blanket, for all $x\in\mathcal X$ and $y\in\mathcal Y$
\[
  \underline{\mathcal R}(y)
  \le \mathcal R_x(y).
\]
Since
$\mathcal R_{\mathrm{BG}}(y)
 = \underline{\mathcal R}(y)/\gamma$, this implies
\[
  \mathcal R_x(y)
  \;\ge\;
  \underline{\mathcal R}(y)
  = \gamma\,\mathcal R_{\mathrm{BG}}(y)
  \qquad\text{for all }x\in\mathcal X,\ y\in\mathcal Y.
\]
Consequently, for all $x\in\mathcal X$ and all $y$ with $\Delta(y)\neq 0$,
\[
  \frac{\Delta(y)^2}{\mathcal R_x(y)}
  \;\le\;
  \frac{\Delta(y)^2}{\gamma\,\mathcal R_{\mathrm{BG}}(y)}.
\]
Integrating over $A(x_1,x_1')$ yields
\[
  \sigma_x^2(x_1,x_1')
  = \int_{A(x_1,x_1')}
      \frac{\Delta(y)^2}{\mathcal R_x(y)}\,dy
  \;\le\;
  \frac{1}{\gamma}
  \int_{A(x_1,x_1')}
      \frac{\Delta(y)^2}{\mathcal R_{\mathrm{BG}}(y)}\,dy
  = \frac{1}{\gamma}\,\sigma_{\mathrm{BG}}^2(x_1,x_1').
\]
Taking the supremum over $x\in\mathcal X$ gives
\[
  \sup_{x\in\mathcal X} \sigma_x^2(x_1,x_1')
  \;\le\;
  \frac{1}{\gamma}\,\sigma_{\mathrm{BG}}^2(x_1,x_1').
\]

By definition,
\[
  \chi_{\mathrm{lo}}(x_1,x_1')
  = \sqrt{\frac{\gamma}{\sigma_{\mathrm{BG}}^2(x_1,x_1')}},
  \qquad
  \chi_{\mathrm{up}}(x_1,x_1')
  = \inf_{x\in\mathcal X}
      \sqrt{\frac{1}{\sigma_x^2(x_1,x_1')}}
  = \frac{1}{\sqrt{\sup_{x\in\mathcal X} \sigma_x^2(x_1,x_1')}}.
\]
Using the inequality on the variances we obtain
\[
  \chi_{\mathrm{up}}(x_1,x_1')
  = \frac{1}{\sqrt{\sup_x \sigma_x^2(x_1,x_1')}}
  \;\ge\;
  \sqrt{\frac{\gamma}{\sigma_{\mathrm{BG}}^2(x_1,x_1')}}
  = \chi_{\mathrm{lo}}(x_1,x_1').
\]

\textbf{Step 4: Characterization of equality.}
We now analyze when equality
$\chi_{\mathrm{up}}(x_1,x_1')=\chi_{\mathrm{lo}}(x_1,x_1')$ holds.
By the definitions above this is equivalent to
\[
  \sup_{x\in\mathcal X} \sigma_x^2(x_1,x_1')
  = \frac{1}{\gamma}\,\sigma_{\mathrm{BG}}^2(x_1,x_1').
\]
From the pointwise inequality
\[
  \frac{\Delta(y)^2}{\mathcal R_x(y)}
  \;\le\;
  \frac{\Delta(y)^2}{\gamma\,\mathcal R_{\mathrm{BG}}(y)},
  \qquad y\in A(x_1,x_1'),
\]
we see that for any fixed $x\in\mathcal X$ we have
\[
  \sigma_x^2(x_1,x_1')
  = \frac{1}{\gamma}\,\sigma_{\mathrm{BG}}^2(x_1,x_1')
  \quad\Longleftrightarrow\quad
  \mathcal R_x(y)
  = \gamma\,\mathcal R_{\mathrm{BG}}(y)
  = \underline{\mathcal R}(y)
  \quad\text{for almost every }y\in A(x_1,x_1').
\]
Indeed, if $\mathcal R_x(y)>\gamma\,\mathcal R_{\mathrm{BG}}(y)$ on a
subset of $A(x_1,x_1')$ of positive measure where $\Delta(y)\neq0$, the
inequality would be strict on that set and hence strict after integration.

Thus equality
$\sup_x \sigma_x^2(x_1,x_1') = \sigma_{\mathrm{BG}}^2(x_1,x_1')/\gamma$
holds if and only if there exists $x^\ast\in\mathcal X$ such that
$\sigma_{x^\ast}^2(x_1,x_1') = \sigma_{\mathrm{BG}}^2(x_1,x_1')/\gamma$,
which in turn is equivalent to
\[
  \mathcal R_{x^\ast}(y)
  = \underline{\mathcal R}(y)
  \quad\text{for almost every }y\in A(x_1,x_1').
\]
Combining this with the previous step proves the lemma.
\end{proof}

\end{document}
